\textbf{Phần 1. Trắc nghiệm lựa chọn}
\setcounterpageref{ex}{0}

\begin{ex} %[3D4N29]
Kết quả của phép nhân $15 \times 4$ là
\choice
{$50$}
{\True $60$}
{$70$}
{$55$}
\loigiai{
$15 \times 4 = 60$.\\
\textbf{Đáp án của câu hỏi là B}
}
\end{ex}

\begin{ex} %[3D4N29]
Kết quả của phép chia $72 : 6$ là
\choice
{$11$}
{$13$}
{$14$}
{\True $12$}
\loigiai{
$72 : 6 = 12$.\\
\textbf{Đáp án của câu hỏi là D}
}
\end{ex}

\begin{ex} %[3D4H29]
Gấp 16 lên 4 lần ta được
\choice
{$54$}
{$74$}
{\True $64$}
{$60$}
\loigiai{
$16 \times 4 = 64$.\\
\textbf{Đáp án của câu hỏi là C}
}
\end{ex}

\begin{ex} %[3D4H29]
Giảm 84 đi 7 lần ta được
\choice
{\True $12$}
{$11$}
{$13$}
{$14$}
\loigiai{
$84 : 7 = 12$.\\
\textbf{Đáp án của câu hỏi là A}
}
\end{ex}

\begin{ex} %[3D4H29]
Tìm số dư trong phép chia $47 : 5$:
\choice
{$1$}
{\True $2$}
{$3$}
{$4$}
\loigiai{
$47 : 5 = 9$ dư 2.\\
\textbf{Đáp án của câu hỏi là B}
}
\end{ex}

\begin{ex} %[3D4H29]
Tính giá trị biểu thức: $18 \times 3 + 26$.
\choice
{\True $80$}
{$78$}
{$82$}
{$84$}
\loigiai{
$18 \times 3 + 26 = 54 + 26 = 80$.\\
\textbf{Đáp án của câu hỏi là A}
}
\end{ex}

\begin{ex} %[3D4V29]
Tính giá trị biểu thức: $95 : 5 - 7$.
\choice
{$11$}
{$13$}
{$14$}
{\True $12$}
\loigiai{
$95 : 5 - 7 = 19 - 7 = 12$.\\
\textbf{Đáp án của câu hỏi là D}
}
\end{ex}

\begin{ex} %[3D4V29]
Một thùng có $24\text{ l}$ dầu, thùng thứ hai có số dầu gấp đôi thùng thứ nhất. Cả hai thùng có số lít dầu là
\choice
{$68\text{ l}$}
{$76\text{ l}$}
{\True $72\text{ l}$}
{$48\text{ l}$}
\loigiai{
Thùng 2: $24 \times 2 = 48\text{ l}$.\\
Tổng: $24 + 48 = 72\text{ l}$.\\
\textbf{Đáp án của câu hỏi là C}
}
\end{ex}

\begin{ex} %[3D4V29]
Có 85 quyển vở chia đều cho 5 tổ. Mỗi tổ có số quyển vở là
\choice
{$15$ quyển}
{$18$ quyển}
{$19$ quyển}
{\True $17$ quyển}
\loigiai{
$85 : 5 = 17$ quyển vở.\\
\textbf{Đáp án của câu hỏi là D}
}
\end{ex}

\begin{ex} %[3D4V29]
Một hình chữ nhật có chiều rộng $12\text{ cm}$, chiều dài gấp 3 lần chiều rộng. Chu vi hình chữ nhật đó là
\choice
{\True $96\text{ cm}$}
{$72\text{ cm}$}
{$48\text{ cm}$}
{$84\text{ cm}$}
\loigiai{
Chiều dài: $12 \times 3 = 36\text{ cm}$.\\
Chu vi: $(36 + 12) \times 2 = 96\text{ cm}$.\\
\textbf{Đáp án của câu hỏi là A}
}
\end{ex}

\begin{ex} %[3D4V29]
Tìm $x$ biết: $x : 6 = 15$.
\choice
{$80$}
{\True $90$}
{$100$}
{$85$}
\loigiai{
$x = 15 \times 6 = 90$.\\
\textbf{Đáp án của câu hỏi là B}
}
\end{ex}

\begin{ex} %[3D4C29]
Một đàn gà có 60 con. Lần thứ nhất bán đi $\dfrac{1}{3}$ số gà, lần thứ hai bán đi $\dfrac{1}{4}$ số gà BAN ĐẦU. Hỏi sau 2 lần bán đàn gà còn lại bao nhiêu con?
\choice
{$20$ con}
{$30$ con}
{\True $25$ con}
{$35$ con}
\loigiai{
Lần 1 bán: $60 : 3 = 20$ con.\\
Lần 2 bán: $60 : 4 = 15$ con.\\
Còn lại: $60 - (20 + 15) = 25$ con.\\
\textbf{Đáp án của câu hỏi là C}
}
\end{ex}

\textbf{Phần 2. Trắc nghiệm đúng sai}
\setcounterpageref{ex}{0}

\begin{ex} %[3D4H29]
Xét tính đúng sai của các phép toán:
\choiceTF
{\True $24 \times 3 = 72$.}
{\True $96 : 4 = 24$.}
{$58 : 6 = 9$ dư $5$.}
{\True $58 : 6 = 9$ dư $4$.}
\loigiai{
\begin{itemize}
    \item $24 \times 3 = 72$: Đúng.
    \item $96 : 4 = 24$: Đúng.
    \item $58 = 6 \times 9 + 4$ (dư 4 $\neq 5$). Ý 3 Sai, Ý 4 Đúng.
\end{itemize}
\begin{itemize}
    \item a) \textbf{Đúng}.
    \item b) \textbf{Đúng}.
    \item c) \textbf{Sai}.
    \item d) \textbf{Đúng}.
\end{itemize}
}
\end{ex}

\begin{ex} %[3D4V29]
Một cửa hàng có $75\text{ kg}$ đường. Buổi sáng bán $\dfrac{1}{5}$ số đường, buổi chiều bán $20\text{ kg}$. Xét tính đúng sai:
\choiceTF
{\True Buổi sáng bán được $15\text{ kg}$ đường.}
{\True Cả ngày bán được $35\text{ kg}$ đường.}
{Cửa hàng còn lại $30\text{ kg}$ đường.}
{\True Cửa hàng còn lại $40\text{ kg}$ đường.}
\loigiai{
\begin{itemize}
    \item Sáng: $75 : 5 = 15\text{ kg}$. Đúng.
    \item Đã bán: $15 + 20 = 35\text{ kg}$. Đúng.
    \item Còn lại: $75 - 35 = 40\text{ kg} \neq 30\text{ kg}$. Ý 3 Sai, Ý 4 Đúng.
\end{itemize}
\begin{itemize}
    \item a) \textbf{Đúng}.
    \item b) \textbf{Đúng}.
    \item c) \textbf{Sai}.
    \item d) \textbf{Đúng}.
\end{itemize}
}
\end{ex}

\textbf{Phần 3. Trả lời ngắn}
\setcounterpageref{ex}{0}

\begin{ex} %[3D4H29]
Tính tích: $17 \times 4$.
\shortans{68}
\loigiai{
\textbf{Đáp số:} $68$\\
$17 \times 4 = 68$.
}
\end{ex}

\begin{ex} %[3D4H29]
Tính thương: $91 : 7$.
\shortans{13}
\loigiai{
\textbf{Đáp số:} $13$\\
$91 : 7 = 13$.
}
\end{ex}

\begin{ex} %[3D4V29]
Gấp 18 lên 5 lần ta được số nào?
\shortans{90}
\loigiai{
\textbf{Đáp số:} $90$\\
$18 \times 5 = 90$.
}
\end{ex}

\begin{ex} %[3D4V29]
Giảm 76 đi 4 lần ta được số nào?
\shortans{19}
\loigiai{
\textbf{Đáp số:} $19$\\
$76 : 4 = 19$.
}
\end{ex}

\begin{ex} %[3D4V29]
Tìm số bị chia trong phép chia có số chia là 5, thương là 16 và số dư là 3.
\shortans{83}
\loigiai{
\textbf{Đáp số:} $83$\\
$5 \times 16 + 3 = 80 + 3 = 83$.
}
\end{ex}

\begin{ex} %[3D4C29]
Một mảnh vườn hình chữ nhật có chu vi $50\text{ m}$, chiều rộng là $10\text{ m}$. Tính chiều dài mảnh vườn theo mét.
\shortans{15}
\loigiai{
\textbf{Đáp số:} $15$\\
Nửa chu vi: $50 : 2 = 25\text{ m}$. Chiều dài: $25 - 10 = 15\text{ m}$.
}
\end{ex}
